\partepigraph{``the speaking of language is part of an activity, or of a form of
life''}{%
Ludwig Wittgenstein, \textit{Philosophical Investigations}}

This part asks how public biomedical resources can support clinical extraction
without hospital records. We study a longitudinal lymphoma electronic case report
form (eCRF), the structured data-collection form of a clinical study, here
comprising 89 fields. The form is filled from a sequence of reports, where each
mention must be assigned to its role in the patient trajectory. No real registry
of this kind is available to us, so we build and evaluate a synthetic version of
the task. The next three chapters construct task supervision from public and
generated text, train an extractor that accepts new field descriptions, and
compare a small biomedical encoder with larger generative models. The complete
eCRF is evaluated on synthetic records, while PARHAF and FrACCO provide narrower
tests on human-written text.

\chapter{Synthetic Data for Task Adaptation}
\label{chap:synthetic}
\section{Introduction}

Clinical text is scarce, but structured medical knowledge is not. France maintains medical ontologies that span tens of thousands of codes (CIM-10 for diagnoses, CCAM for procedures, ATC for medications), and PubMed Central contains millions of clinical case descriptions embedded in scientific articles. Both are publicly available, both are richly structured, and neither can be used directly to train a language model. The question is whether we can convert them into synthetic training data that is closer to what clinical applications actually need.

We already saw one form of generated data in \Cref{chap:ontobook}, where medical ontologies became synthetic textbooks for pretraining. This chapter generates data for the extraction task itself, in two ways, and neither touches a real patient record. The first annotates public biomedical text with a large language model, producing the labelled examples that train the extraction model of \Cref{chap:architectures}. The second rewrites public case descriptions into the style of hospital notes, to approximate a register that no public corpus offers. Both take a public source, hand it to a language model under a strict prompt, and get back training material that real clinical data could not supply.


\section{Annotating Biomedical Text for Extraction}
\label{sec:gliner-data}

\begin{figure}[t]
\centering
\providecommand{\ftitle}[1]{{\normalsize\bfseries\sffamily\color{ThesisInk}#1}}
\providecommand{\etype}[1]{{\color{ThesisPrimaryDark}#1}}
\resizebox{\textwidth}{!}{%
\begin{tikzpicture}[
  font=\small\sffamily,
  outbox/.style={draw, rounded corners=4pt, line width=0.8pt,
                 text width=7.4cm, inner sep=12pt, align=left, font=\small},
]
% Input text block (left)
\node[draw=ThesisNeutral, fill=ThesisPaper, rounded corners=4pt, line width=0.8pt,
      text width=3.7cm, inner sep=11pt, align=justify, font=\small,
      text=ThesisInk] (input) at (0,0)
  {\itshape Mrs Diallo, 64, diffuse large B-cell lymphoma diagnosed 12/03/2019.
   R-CHOP started 02/04/2019, 6 cycles. Relapse 15/01/2021. Death 03/02/2021.
   LDH 480 IU/L at diagnosis.};

% Three output blocks (right), stacked with generous spacing
\node[outbox, draw=ThesisPrimaryDark, fill=ThesisPrimary!10,
      anchor=west] (ent) at (5.6,4.2)
  {\ftitle{Entities}\par\vspace{5pt}
   {\ttfamily\etype{Disease}\,$\rightarrow$\,diffuse large B-cell lymphoma\par
    \etype{Drug}\,$\rightarrow$\,R-CHOP\quad
    \etype{Examination}\,$\rightarrow$\,LDH\par
    \etype{date}\,$\rightarrow$\,03/02/2021\quad
    \etype{measure}\,$\rightarrow$\,480 IU/L}};

\node[outbox, draw=ThesisSecondaryDark, fill=ThesisSecondary!12,
      anchor=west] (rel) at (5.6,0)
  {\ftitle{Relations}\par\vspace{5pt}
   {\ttfamily treats:\, R-CHOP\,$\rightarrow$\,lymphoma}};

\node[outbox, draw=ThesisTertiaryDark, fill=ThesisTertiary!14,
      anchor=west] (str) at (5.6,-4)
  {\ftitle{Structure \& classification}\par\vspace{5pt}
   {\ttfamily timeline \{diagnosis date, treatment start,\par
    \hphantom{timeline \{}relapse, death\}\par
    ICD-10 chapter: tumours\par
    document type: follow-up note}};

% Fan-out arrows through an LLM label
\coordinate (hub) at (4.2,0);
\draw[-{Latex[length=2.6mm]}, line width=1.1pt, ThesisInk!55]
   (input.east) -- node[above, font=\small\bfseries, color=ThesisInk]{LLM} (hub);
\draw[-{Latex[length=2.6mm]}, line width=1pt, ThesisInk!45] (hub) to[out=55,in=180] (ent.west);
\draw[-{Latex[length=2.6mm]}, line width=1pt, ThesisInk!45] (hub) -- (rel.west);
\draw[-{Latex[length=2.6mm]}, line width=1pt, ThesisInk!45] (hub) to[out=-55,in=180] (str.west);
\end{tikzpicture}%
}
\caption{A large language model turns one passage into typed entities, relations
between them, and a document structure with its classifications. The full
annotation prompts are given in \Cref{app:gliner-prompts}.}
\label{fig:annotation-example}
\end{figure}

Training an extractor by fine-tuning needs labelled examples, and as this part's introduction noted, those are out of reach for French clinical text. We take a different route. Instead of paying clinicians to annotate real notes, we annotate public biomedical text automatically with a large language model, and use the result to train a smaller extraction model. This is a form of distillation: a large general model that is too heavy to run at scale produces the labels, and a small model learns to reproduce them. Everything we annotate is public text, so none of the resulting data carries the restrictions of a patient record.

The annotator is \texttt{Qwen3-235B-A22B-Instruct}, a 235-billion-parameter instruction-tuned model~\citep{yang2025qwen3}, run with FP8 quantisation on four NVIDIA H100 GPUs through vLLM~\citep{kwon2023vllm}. For each document, the model returns a single JSON object whose fields are fixed in advance by a schema, so every annotation has the same shape and parses without error. We turn off the model's step-by-step reasoning mode and sample at a low temperature of $0.6$, which keeps the output close to the schema. The same model later serves as the judge that scores our extractors (\Cref{chap:evaluation}), so we check there that it does not simply favour its own annotations.

The schema asks for more than a list of entities. For each document the model first lists candidate spans, meaning every noun phrase that could be a medical entity, and then builds four kinds of annotation on top of them. It marks \emph{entities} with open-vocabulary types at four levels of granularity, from a single reusable word such as \textit{Maladie} (disease) up to a full compositional phrase, and the same span may be typed at several levels at once. It assigns two or three \emph{classifications} to the document, such as its specialty or its document type. It fills zero to two \emph{structures}, small field-value schemas such as a dosage or a patient profile. And it records \emph{relations} between entities, such as a drug that treats a disease. The model also returns a short list of \emph{hard negatives}: plausible types that are absent from the text, which teach the extractor what not to find. \Cref{fig:annotation-example} shows one annotated example.

We then ground every annotation against the source text. A mention is kept only when it is an exact substring of the document, which removes spans the model paraphrased or invented. We also drop a type when it is merely a lexical variant of its own mention, measured by edit distance, so the model cannot take the shortcut of repeating a word as its own type. The head and tail of every relation must likewise be exact substrings. Documents where no annotation survives are not discarded: a document with text but no entity is kept as a negative example, which is what teaches the extractor to stay silent on non-medical text.

The text we annotate is drawn from three kinds of public source, listed in \Cref{tab:gliner-sources}. Most of it is French biomedical literature from the corpus of \Cref{chap:mcbio}, kept above a quality threshold. A large share is case-centred: articles that contain patient cases, and French translations of PubMed Central case reports. A smaller part comes from passages generated from the CIM-10, CCAM, and ATC nomenclatures, which bring in coding terminology. Finally, non-medical and low-quality documents are added on purpose as negatives. The final dataset holds about $94{,}000$ annotated documents, with $662{,}000$ entities, $1.98$~million mentions, and $592{,}000$ relations; roughly a quarter of the documents are negatives.

\begin{table}[t]
  \centering
  \caption{Public sources sampled for annotation by the language model to train
  the extraction model of \Cref{chap:architectures}. Of the $118{,}000$ sampled
  documents, about $94{,}000$ yield a valid annotation after grounding and
  filtering.}
  \label{tab:gliner-sources}
  \small
  \setlength{\tabcolsep}{12pt}
  \renewcommand{\arraystretch}{1.25}
  \begin{tabular}{@{}ll S[table-format=6.0, group-separator={,}, group-minimum-digits=4]@{}}
    \toprule
    Source & What it provides & {Sampled} \\
    \midrule
    \tablegroup{3}{Biomedical literature (\Cref{chap:mcbio})}
    HAL                 & French biomedical articles             & 26000 \\
    ISTEX               & French biomedical articles             & 26000 \\
    \tablegroup{3}{Patient cases}
    Articles with cases & passages describing real patient cases  & 18000 \\
    PMC-Patients-FR     & translated PubMed Central case reports  & 14000 \\
    \tablegroup{3}{Ontology terminology}
    CIM-10 / CCAM / ATC & passages from the coding nomenclatures  & 9000 \\
    \tablegroup{3}{Negatives}
    Non-medical / low quality & documents with no entity to find  & 25000 \\
    \midrule
    Documents sampled    &                                        & 118000 \\
    Valid annotations    & after grounding and filtering          & 94000 \\
    \bottomrule
  \end{tabular}
\end{table}



\section{Synthetic Clinical Reports}
\label{sec:synthetic-reports}

Public case descriptions still read like published cases, not like the terse notes a hospital produces. The second use of generation closes part of that gap. We rewrite a public case into the register of a hospital note, and a deterministic degradation step adds the abbreviations, missing accents, and dictation slips that make real notes hard to read, so the extractor meets that register during training as well as on real notes.


\section{Synthesizing Longitudinal Patient Dossiers}
\label{sec:lymphoma-dossiers}

A single note captures one moment. A patient's record is a story that unfolds over years and across many documents, and the clinical questions we care about live in that story rather than in any one note. This thesis is part of OncoLab, a project that annotates the dossiers of lymphoma patients, all of them relapsed, against a form of about a hundred variables covering demographics, disease history, treatment lines, imaging, and outcome. Real dossiers cannot be used to develop and test that annotation, for the privacy reasons that run through this thesis, so the project needs a synthetic but realistic set of French dossiers to build the form and train automatic extraction before any real record is touched.

A first attempt seeded each synthetic dossier from a single short case description, and it did not work. The reason is worth stating, because it shaped everything that followed. The seed text was about a thousand characters and described one moment, while a dossier spans years and a dozen documents. Asked to bridge that distance, the model had to invent almost everything, and it did so by imitation, producing dossiers that looked plausible but did not hold together. Clinicians who reviewed the output found the revealing gaps: there were almost no follow-up reports and almost no relapse consultations, although every patient in the target cohort has relapsed; cases with both a first and a second treatment line were missing, so adverse events beyond the first cycle could not be annotated; and the multidisciplinary-meeting reports had the wrong shape.

The fix was not to add a rule for each defect. Patching the prompt with one constraint per problem makes the model produce clone-like patients and loses the variety that made it useful. The better question is what lets the model do well on its own, and the answer was a better source and a decomposed task. We seed instead from real longitudinal case reports, translated into French, in which relapse, second-line treatment, adverse events, and death are already present in the text. From the $167{,}034$ public cases we keep the richest, giving more weight to those that mention a relapse and a second treatment line, down to about $2{,}800$ adult cases.

The generation then splits into stages and respects time. The model first writes a coherent patient timeline as structured data, and a short reflexive pass repairs any inconsistency in that timeline rather than rejecting it. Only then does it write the documents, one at a time and in chronological order: each document is generated seeing the earlier ones but never the later ones, so the dossier accumulates the way a real record does and no document can quote a result from its own future. A final pass extracts the structured fields. Because a field such as a treatment date takes several values over a dossier, each field is stored as a list of values, each tied to the document it came from, which keeps the longitudinal structure and lets every value be traced back to its source.

The redesign produced about three thousand synthetic dossiers, around eight documents each, with a valid timeline and a complete annotation in almost every case. Second-line treatment now appears in about a third of the dossiers, against roughly a fifth before, which is the gap the clinicians had flagged. The work is still in progress: the most recent run covers $2{,}050$ of its $2{,}774$ planned cases, and the automatic quality check measures format coherence rather than clinical accuracy, so it should not be read as a correctness score. A held-out part of this set becomes the hardest evaluation in the thesis, the one we use in \Cref{chap:evaluation}.


\section*{Conclusion}

Across these uses the recipe is the same: take a public source, constrain a language model with a strict prompt, and keep only what holds together. The data this produces is real and large, but it settles nothing on its own. It only matters if a model can turn it into clinical extraction under a label set that is open and keeps changing, the setting that defeats a fixed classifier.

\vspace{2em}
The data now exists. Whether a model can use it well is the question of the next chapter, \Cref{chap:architectures}, which builds the extractor trained on the material generated here.


\chapter{Architectures for Low-Resource Extraction}
\label{chap:architectures}
\section{Introduction}

The previous chapter produced training data without a fixed label set, and the introduction to this part explained why a fixed label set is the wrong fit for clinical extraction: the types we want change from one study to the next, and most will never have been seen in training. A model with a fixed output layer cannot produce a type it was not trained on. This chapter builds a model that has no fixed output layer at all. Instead of choosing among a closed list of labels, it compares a span of text against a short description of the type we are looking for and decides whether they match. The type is given at the moment of use, so one model can extract a disease, a dosage, or a treatment response without being retrained for each.

\section{Open-Vocabulary Extraction by Type Description}

This is the idea behind GLiNER~\citep{zaratiana2024gliner}: rather than learning one output per label, the model learns to compare. It reads the document together with the descriptions of the types we are looking for, and scores how well each span of text matches each type. Because the types are supplied as text at inference time, the set of types is open: a user can ask for a type the model never saw in training, and the model will still try to find it.

Our extractor, \emph{ModernCamemBERT-bio-gliner-base}, follows the GLiNER2 design~\citep{zaratiana2025gliner2}. A single encoder reads the document and the candidate type descriptions in one pass, and a light span-scoring head matches every span against every type. The same machinery also handles the other tasks the training data carries: classifying the document, filling structured fields, and extracting relations between spans.

We release the model in two configurations that differ only in their training-data mix. The first is tuned for entity extraction and is the stronger general extractor: it finds types it was never trained on more accurately than any baseline we compare against in \Cref{chap:evaluation}. The second is multi-task: alongside entities it classifies the document, fills structures, and extracts relations, and it does these better than the first, at some cost to entity extraction. The trade-off between the two is itself one of the findings of this part, and we return to it in \Cref{chap:evaluation}.

\section{The MedEmbed Backbone}

The matching is only as good as the encoder that places type descriptions and spans in a shared space. We train our own, \emph{MedEmbed}, a French biomedical sentence encoder. It starts from ModernCamemBERT-bio, the modern French biomedical encoder developed in \Cref{chap:objectives,chap:ontobook}, and is trained with a contrastive objective: a type description and a matching span are pulled together, and mismatched pairs are pushed apart. The model has 149 million parameters and reads up to 8192 tokens, so a long description and a long passage fit in one pass.

The training pairs that matter most come from medical terminologies. The CIM-10, CCAM, and ATC nomenclatures are turned into description-span pairs by the ontology work of \Cref{chap:ontobook}, which we do not repeat here.

On OntoBench-FR, a retrieval benchmark built from the French terminologies and from real clinical cases, MedEmbed reaches 58.1\% average accuracy and beats encoders several times its size: 48.9\% for mE5-large-instruct, 47.2\% for BGE-M3, and 44.9\% for Solon.\footnote{This is the average weighted by the number of test items per benchmark; an unweighted average gives a slightly higher figure on a different scale, and the two should not be compared directly.} An ablation confirms where the strength comes from: the terminology pairs are the one source that matters, and training on them alone is enough to match the full mixture, while removing them costs more than ten points. The starting encoder matters too: ModernCamemBERT-bio gives a clear gain over a plainly pretrained encoder on the novel types that generalisation depends on.

\section{Training}

We train the extractor on the synthetic annotations of \Cref{chap:synthetic}, with no real clinical labels at any point. Both configurations use the same recipe and run in a few hours on a single GPU. The result is a small open model that extracts types defined at inference time, trained entirely on public and synthetic data.

\section*{Conclusion}

A model that can extract a type it was never trained on raises a problem the moment we try to measure it. The usual way to score extraction is to compare against a fixed human annotation, but a fixed annotation only covers the types the annotators were asked to mark, which is exactly the closed world this model was built to escape. Measuring it fairly needs a different kind of evaluation, and that is the subject of the next chapter.


\chapter{Evaluating Open-Vocabulary Extraction}
\label{chap:evaluation}
\section{Introduction}

The extractor of the previous chapter can be asked for a type it never saw. That is the point of it, and it is also what makes it hard to score. The usual way to measure extraction compares a model's output against a fixed set of human annotations. But those annotations cover only the types the annotators were told to mark, under one set of conventions about where a span begins and ends and what counts as an entity. A model rewarded for matching those conventions is not necessarily better at finding what a new study cares about; it may have learned the house style of one annotation effort. To see whether our model generalises, we need an evaluation that does not assume a fixed answer key.

\section{A Reference-Free Evaluation with a Language-Model Judge}

We evaluate without a gold standard. We take 400 clinical case documents the models were never trained on, and a bank of twenty types: ten common clinical types, such as disease, drug, and symptom, and ten novel types that no benchmark uses, built to be compositional or cross-domain, such as an adverse effect, an implantable device, or an environmental risk factor. Each model runs twice, once for the common types and once for the novel ones, so the two can be measured apart. We pool the spans every model proposes, add candidates of the judge's own, remove duplicates, and hide which system produced which span. A large language model then judges each span against its type description and labels it correct, wrong type, wrong boundary, or not an entity. We score with mean average precision, which does not depend on a decision threshold, and we summarise generalisation by the gap between a model's score on common types and its score on novel ones.

To keep the judge honest we seed the pool with controls: real clinical entities that should be accepted, and function words mislabelled as diseases that should be rejected. We keep only the runs where the judge handles these correctly.

\section{Benchmark Winners Do Not Generalise}

On the common types the models are close. The novel types separate them. Our entity-tuned configuration reaches a mean average precision of $0.476$ on novel types, the best of any model we test. The baselines collapse: a strong single-encoder model tuned to a standard benchmark scores well on common types but never makes a confident prediction on a novel one, and a published biomedical extractor falls to $0.121$ on novel types. The gap between common and novel performance, which measures how far a model has overfit to familiar types, is small for our models and large for the baselines.

Ranked on a standard benchmark, the same models come out in almost the opposite order: the convention-tuned baseline first, our multi-task configuration last. A model can top the benchmark and still be the weakest at finding new types, because the benchmark measures agreement with one annotation convention rather than the ability to generalise. This is the central result of the part. A lower benchmark score is what we should expect from a model that has not overfit to a single convention, and it is the entity-tuned configuration, weaker on the benchmark, that generalises best.

\section{Is the Judge Trustworthy?}

Using a language model to grade extractions invites two worries: that the judge is unreliable, and that it favours models like itself. We check both. Against human annotations on clinical types where the gold is unambiguous, the judge agrees at a Cohen's kappa of $0.532$, rising to $0.797$ on the cleanest set. Two judges from different model families agree with each other at $0.728$, so the result does not rest on a single judge. Re-scoring everything with a second, independent medical model leaves the rankings in place and slightly sharper, which rules out self-preference. And where the judge and the human gold disagree, $85\%$ of the disagreements are cases where the human annotation followed a convention the judge reasonably declined to follow, not cases where the judge was wrong.

The judge is conservative: it under-credits correct spans, so the absolute scores understate real performance while the rankings between models stay stable. One limit remains. The human check covers only the common types, so the headline claim about novel types has no direct human anchor yet, and a small human study of novel-type judgements is the most useful next step.

\section{The OncoLab Lymphoma Cohort}

The last test is the one that matters most, because it puts the whole thesis to work on a real clinical problem. The OncoLab cohort asks for $89$ fields of a lymphoma annotation form, and we evaluate on the held-out synthetic dossiers built in \Cref{chap:synthetic}, scoring each field by how well the model locates its value in the text. Every earlier part of the thesis meets here: the public corpora that supplied the seeds (\Cref{chap:collecting,chap:quality,chap:mcbio}), the modern French biomedical encoder (\Cref{chap:objectives,chap:ontobook}), the open-vocabulary architecture (\Cref{chap:architectures}), and the synthetic dossiers themselves (\Cref{chap:synthetic}).

Off the shelf, no model does well. The best zero-shot result is a micro-averaged F-score of $0.174$, and most models sit between $0.10$ and $0.15$. The convention-tuned baseline is precise but barely fires, finding only a small fraction of the fields. This is a hard, long, multi-document task, and reading it cold is beyond every model we try.

After fine-tuning on the synthetic dossiers, the picture changes. Every model rises to about $0.31$ to $0.33$, and the differences between them nearly vanish. The synthetic data closes most of the distance, and it closes it for all of them: the benchmark champion and the generalist extractor reach the same ceiling. The scores are still low in absolute terms, which measures how hard the task is rather than showing it solved. But the shape of the result is the payoff of this work. Built only from public and synthetic data, a small open model reaches the same level on a real oncology task as a model tuned on a standard benchmark, without ever seeing a real clinical note.

\section*{Conclusion}

This is where the parts of the thesis come together. A public biomedical corpus, an encoder trained on it, a training objective with a dose of structured knowledge, an open-vocabulary architecture, and data that is generated rather than collected, all converge on one relapsed-lymphoma cohort and let a model that never saw a real clinical note reach a usable level of extraction. The conclusion of the thesis steps back from these numbers to ask what they say about building clinical models from public data, and where the approach still falls short.

